\subsection{Periodic closure of the single-wall spaces}

% Source-written declarations; hosted elaboration and declaration checks are pending.
\begin{theorem}[Cyclic boundary trace coefficients]
    \label{thm:uncle_cyclic_boundary_coefficients}
    \lean{MPSTensor.trace_uncleBlockMatrix_mul,
        MPSTensor.eq_of_uncleBoundary_cyclic_pairing,
        MPSTensor.exists_scalar_blocks_of_uncleBoundary_trace_cyclic}
    \notready
    \uses{def:two_block_uncle_tensor}
    On four blocks write
    \begin{align}
        (A,R,L,B)\star(C,S,T,E)
        =(AC,AS+RE,LC+BT,BE).
        \notag
    \end{align}
    If two boundary matrices $X,Y$ satisfy
    $\operatorname{tr}((M\star N)X)=\operatorname{tr}((N\star M)Y)$
    for every $M,N$, then $X=Y$. If the pairing with $X$ is cyclic,
    then $X=\operatorname{diag}(\alpha I,\beta I)$ for some scalars
    $\alpha,\beta$.
\end{theorem}

\begin{proof}
    Taking $M=I$ gives $X=Y$ by nondegeneracy of the ordinary trace
    pairing. Pair a sector identity with an off-diagonal block: one
    product is that block and the reversed product is zero. Trace
    nondegeneracy therefore kills both off-diagonal boundary blocks.
    On either diagonal sector, cyclicity says
    $\operatorname{tr}(PQX)=\operatorname{tr}(QPX)$. Trace
    nondegeneracy makes $X$ commute with every sector matrix, so each
    diagonal block is scalar.
\end{proof}

\begin{theorem}[Scalar boundaries are the periodic sector vectors]
    \label{thm:uncle_scalar_boundaries}
    \lean{MPSTensor.uncleBoundaryMap_scalar_blocks}
    \notready
    \uses{def:two_block_uncle_tensor}
    For the actual single-wall boundary map $\Gamma_N$, the boundary
    $\operatorname{diag}(\alpha I,\beta I)$ gives
    $\Gamma_N(\operatorname{diag}(\alpha I,\beta I))
        =\alpha\ket{M(A)}+\beta\ket{M(B)}$.
\end{theorem}

\begin{proof}
    The diagonal blocks of the single-wall coefficient word are the
    pure $A$ and $B$ words. Scalar diagonal boundaries remove the two
    off-diagonal contributions, leaving their two traces.
\end{proof}

\begin{theorem}[Periodic uncle boundary closure]
    \label{thm:uncle_periodic_boundary_closure}
    \lean{MPSTensor.uncleBoundarySpace_inf_cyclicTranslate}
    \notready
    \uses{def:two_block_uncle_tensor, def:parent_cyclic_translate_state}
    Let the full perturbation letter
    $\left(\begin{smallmatrix}A&R\\L&B\end{smallmatrix}\right)$ be
    injective. For every $N\ge2$, let $S_N$ be the actual sum of the
    two pure-sector and two single-wall boundary spaces, and let $T$
    be the one-site cyclic translation. Then
    $S_N\cap T(S_N)=\operatorname{span}\{\ket{M(A)},\ket{M(B)}\}$.
    This is the closure lemma in~\cite{FernandezGonzalez2015Uncle},
    used after the two-site uncle interaction.
\end{theorem}

\begin{proof}
    \uses{thm:uncle_cyclic_boundary_coefficients,
        thm:uncle_scalar_boundaries}
    Represent a vector in the intersection by boundaries $X$ and $Y$
    at the two cuts. A word with first letter $i$ and remaining
    coefficient $W$ gives
    $\operatorname{tr}((C_i\star W)X)
        =\operatorname{tr}((W\star C_i)Y)$.
    Full perturbation injectivity makes both the one-letter
    coefficients and the positive-length remaining words span all
    coefficient matrices. Extend the identity linearly in each
    variable. The cyclic boundary theorem makes $X=Y$ scalar on
    each diagonal sector, giving the asserted span. Conversely,
    each periodic sector vector belongs to $S_N$ and is fixed by
    cyclic translation of its trace.
\end{proof}
